% ============================================================
%  附录C 双时格林函数的运动方程（原书印刷页 543–546）
%  底本：高等教育出版社《固体理论》第二版
%  （PDF 页 = 印刷页 + 20，原书 PDF 页 563–566）
%  公式编号自检：全附录 (C.1)–(C.21) 连续 21 个，无缺号。
%    任务单 tags=19 系 (C.3) 编号嵌于 OCR 乱码、(C.19) 编号游离
%    于两式之间未被计入（原书 p563、p545 均在号）。
%
%  已修正的 OCR/原书错误（除全角句点转"."、删 ~ 间隔符、
%  拆开行内公式与中文等常规项外）：
%  1. (C.1)(C.3)(C.9)(C.10) 整块 OCR 乱码按 p563–564 复原；
%     (C.10) 跨页两行（p564 页末 + p565 页首）已并为一式。
%  2. (C.8) 无穷序列 OCR 末行杂符（\tilde{A}_{n+1} 之衍符、
%     重复的 \vdots 等）按 p564 清理复原。
%  3. 原书 p545 (C.17) 印"Σ_r(k,ω)≡Σ(k,ω−iη)"，按 (C.15) 分母
%     Σ(k,ω+iη) 及推迟自能应在 上半复平面解析 判定 −iη 为
%     +iη 之误，改 Σ(k,ω+iη)。
%  4. (C.19) 原书为独立编号公式（"或"之前），OCR 把编号游离成
%     孤行，已归还本式。
%  5. 体例：< <A|B> > 统一 \langle\langle…\rangle\rangle；虚单位
%     i、指数底 e 用正体；宗量 k 用粗体（照原书）。
%  遗留问题：
%  · (C.17) 之更正属原书排印疑误，已按上下文订正，见上第 3 条。
% ============================================================
\chapter{双时格林函数的运动方程}

利用附录 A 中定义的、包括了 $G_{\mathrm{r}}(\omega)$ 与 $G_{\mathrm{a}}(\omega)$ 的全格林函数 $G(\omega)=\langle\langle A|B\rangle\rangle_{\omega}$，作傅里叶变换可相应地定义全双时格林函数
\begin{equation}
\begin{aligned}
G(t)&\equiv\frac{1}{2\pi}\int_{-\infty}^{\infty}
\mathrm{d}\omega\,\mathrm{e}^{-\mathrm{i}\omega t}G(\omega)
\equiv\langle\langle A(t);B\rangle\rangle\\
&\equiv\left\{
\begin{aligned}
&-\mathrm{i}\theta(t)\langle[A(t),B]_{\pm}\rangle
=G_{\mathrm{r}}(t)\\
&\mathrm{i}\theta(-t)\langle[A(t),B]_{\pm}\rangle
=G_{\mathrm{a}}(t)
\end{aligned}\right.
\end{aligned}
\end{equation}
显然，$G(t>0)$ 代表推迟格林函数，而 $G(t<0)$ 为超前格林函数. $A(t)$ 为海森伯绘景中算符 $\mathrm{e}^{\mathrm{i}Ht}A\mathrm{e}^{-\mathrm{i}Ht}$，这里取 $\hbar=1$ 单位. 从上式出发，对 $t$ 求导，并利用海森伯运动方程
\begin{equation}
\mathrm{i}\frac{\mathrm{d}A(t)}{\mathrm{d}t}
=[A(t),H]=A(t)H-HA(t)
\end{equation}
求出格林函数的运动方程
\begin{equation}
\begin{aligned}
\mathrm{i}\frac{\mathrm{d}}{\mathrm{d}t}
\langle\langle A(t);B\rangle\rangle
&=\delta(t)\langle[A(t),B]_{\pm}\rangle
\mp\mathrm{i}\theta(\pm t)
\langle[[A(t),H],B]_{\pm}\rangle\\
&=\delta(t)\langle[A,B]_{\pm}\rangle
+\langle\langle[A(t),H];B\rangle\rangle
\end{aligned}
\end{equation}
再利用
\begin{equation}
\delta(t)=\frac{1}{2\pi}\int_{-\infty}^{\infty}
\mathrm{d}\omega\,\mathrm{e}^{-\mathrm{i}\omega t}
\end{equation}
\begin{equation}
\begin{aligned}
\mathrm{i}\dot{G}(t)&=\frac{1}{2\pi}\int_{-\infty}^{\infty}
\mathrm{d}\omega\,\omega G(\omega)\mathrm{e}^{-\mathrm{i}\omega t}\\
&=\frac{1}{2\pi}\int_{-\infty}^{\infty}\mathrm{d}\omega\,\omega
\langle\langle A|B\rangle\rangle_{\omega}\mathrm{e}^{-\mathrm{i}\omega t}
\end{aligned}
\end{equation}
而在式(C.3)中出现的高阶格林函数也可作傅里叶变换，用复平面内的 $\omega$ 分量表示为
\begin{equation}
\langle\langle[A(t),H];B\rangle\rangle=\frac{1}{2\pi}
\int_{-\infty}^{\infty}\mathrm{d}\omega\,
\langle\langle[A,H]|B\rangle\rangle_{\omega}
\mathrm{e}^{-\mathrm{i}\omega t}
\end{equation}
将式(C.4)—(C.6)代入式(C.3)，求得在复平面 $\omega$ 上的格林函数运动方程
\begin{equation}
\omega\langle\langle A|B\rangle\rangle_{\omega}
=\langle[A,B]_{\pm}\rangle
+\langle\langle[A,H]|B\rangle\rangle_{\omega}
\end{equation}
注意，上式右边的 $\langle\langle[A,H]|B\rangle\rangle_{\omega}$ 是高阶格林函数，它的运动方程仍为式(C.7)，只不过其中需用 $A_{1}=[A,H]$ 代替 $A$. 按此类推，可认为格林函数方程式(C.7)是适用于各阶格林函数的一个通式，并构成一组耦合格林函数方程的无穷序列：
\begin{equation}
\left.
\begin{aligned}
&\omega\langle\langle A|B\rangle\rangle_{\omega}
=\langle[A,B]_{\pm}\rangle
+\langle\langle A_{1}|B\rangle\rangle_{\omega}
\qquad A_{1}\equiv[A,H]\\
&\omega\langle\langle A_{1}|B\rangle\rangle_{\omega}
=\langle[A_{1},B]_{\pm}\rangle
+\langle\langle A_{2}|B\rangle\rangle_{\omega}
\qquad A_{2}\equiv[A_{1},H]\\
&\qquad\vdots\\
&\omega\langle\langle A_{n}|B\rangle\rangle_{\omega}
=\langle[A_{n},B]_{\pm}\rangle
+\langle\langle A_{n+1}|B\rangle\rangle_{\omega}
\qquad A_{n+1}\equiv[A_{n},H]\\
&\qquad\vdots
\end{aligned}
\right\}
\end{equation}
其中 $A_{n+1}\equiv\left[\cdots\left[A,H\right],\cdots,H\right]$ 代表 $H$ 的 $(n+1)$ 阶泊松括号. 显然随着 $n$ 增大，将涉及更高阶的格林函数 $\langle\langle A_{n+1}|B\rangle\rangle_{\omega}$，这是一组无穷的联立方程系统，它是不封闭的. 因为若要严格求解 $\langle\langle A|B\rangle\rangle_{\omega}$，就必须得知 $\langle\langle A_{1}|B\rangle\rangle_{\omega}$，而后者又是由 $\langle\langle A_{2}|B\rangle\rangle_{\omega}$，$\cdots$，$\langle\langle A_{n}|B\rangle\rangle_{\omega}$,$\cdots$ 以至无穷高阶的格林函数决定的，这是由 $H$ 中的相互作用所引起的困难，从一个侧面说明了多体问题一般讲是不能严格求解的现实. 在具体求解时，往往需要对高阶格林函数作某种近似（称为切断近似），使式(C.8)在有限阶封闭，从而可用代数方法解出近似的 $\langle\langle A|B\rangle\rangle_{\omega}$. 在某些已取特殊极限或作了近似的简化模型中，方程组自动在某阶封闭，这时所求出的格林函数常称为该简化模型的严格解，而这些简化模型也相应的称为严格可解模型. 例如在原子极限下的赫伯德模型和在自洽场近似下的安德逊模型等，我们将在第十一章§1与§2中作专门的讨论.

对于双时格林函数 $\langle\langle A|B\rangle\rangle_{\omega}$ 也可以按照戴逊方案引入自能函数来处理上述耦合方程的无穷序列. 例如对于 $H=H_{0}+H_{\mathrm{int}}$ 的多体系统，将式(C.7)右边 $\langle\langle[A,H]|B\rangle\rangle_{\omega}$ 中的相互作用部分 $\langle\langle[A,H_{\mathrm{int}}]|B\rangle\rangle_{\omega}$ 在形式上写成 $\Sigma(\omega)\langle\langle A|B\rangle\rangle_{\omega}$，就可以得到格林函数的戴逊方程，其中对 $\langle\langle A|B\rangle\rangle_{\omega}$ 的互作用修正已完全概括到自能函数 $\Sigma(\omega)$ 之中了. 我们将以相互作用的电子系统为例作简要的介绍.

设互作用电子系统的哈密顿量为
\begin{equation}
\left.
\begin{aligned}
&H=H_{0}+H_{\mathrm{int}}\\
&H_{0}=\sum_{k,\sigma}E_{k}C_{k\sigma}^{+}C_{k\sigma}
\end{aligned}
\right\}
\end{equation}
而 $H_{\mathrm{int}}$ 为含四算子的二体互作用部分，利用格林函数运动方程的通式(C.7)，单电子格林函数满足方程
\begin{equation}
\begin{aligned}
\omega\langle\langle C_{k\sigma}|C_{k\sigma}^{+}\rangle\rangle_{\omega}
&=\langle[C_{k\sigma},C_{k\sigma}^{+}]_{+}\rangle
+\langle\langle[C_{k\sigma},H_{0}]|C_{k\sigma}^{+}\rangle\rangle_{\omega}\\
&\qquad+\langle\langle[C_{k\sigma},H_{\mathrm{int}}]
|C_{k\sigma}^{+}\rangle\rangle_{\omega}
\end{aligned}
\end{equation}
其中等式右边一、二两项可直接算出为
\begin{equation}
[C_{k\sigma},C_{k\sigma}^{+}]_{+}=1\text{，}\qquad
[C_{k\sigma},H_{0}]=E_{k}C_{k\sigma}
\end{equation}
得到方程
\begin{equation}
(\omega-E_{k})\langle\langle C_{k\sigma}
|C_{k\sigma}^{+}\rangle\rangle_{\omega}
=1+\langle\langle[C_{k\sigma},H_{\mathrm{int}}]
|C_{k\sigma}^{+}\rangle\rangle_{\omega}
\end{equation}
由于 $H_{\mathrm{int}}$ 中含有 4 个 C 算子，$[C_{k\sigma},H_{\mathrm{int}}]$ 必含 3 个 C 算子，因此 $\langle\langle[C_{k\sigma},H_{\mathrm{int}}]|C_{k\sigma}^{+}\rangle\rangle_{\omega}$ 是含 4 个算子的双粒子格林函数：根据式(C.8)的分析可知，它又涉及三粒子以至更多粒子的格林函数，这样在 $\langle\langle[C_{k\sigma},H_{\mathrm{int}}]|C_{k\sigma}^{+}\rangle\rangle_{\omega}$ 的严格解中应当包括了相互作用的全部信息，并可以按照下列方式引进自能函数 $\Sigma(\boldsymbol{k},\omega)$ 来概括相互作用对 $\langle\langle C_{k\sigma}|C_{k\sigma}^{+}\rangle\rangle_{\omega}$ 的全部修正
\begin{equation}
\langle\langle[C_{k\sigma},H_{\mathrm{int}}]
|C_{k\sigma}^{+}\rangle\rangle_{\omega}
=\Sigma(\boldsymbol{k},\omega)
\langle\langle C_{k\sigma}|C_{k\sigma}^{+}\rangle\rangle_{\omega}
\end{equation}
其中 $\omega$ 为复平面上的量，式(C.13)使格林函数运动方程在形式上退耦而写成单粒子格林函数方程
\begin{equation}
[\omega-E_{k}-\Sigma(\boldsymbol{k},\omega)]
\langle\langle C_{k\sigma}|C_{k\sigma}^{+}\rangle\rangle_{\omega}=1
\end{equation}
而将所有的复杂修正都概括到未知的自能函数之中. 这时双时推迟格林函数可简单的表示为
\begin{equation}
G_{\mathrm{r}}(\boldsymbol{k},\omega)
=\langle\langle C_{k\sigma}|C_{k\sigma}^{+}\rangle\rangle_{\omega+\mathrm{i}\eta}
=\frac{1}{\omega-E_{k}+\mathrm{i}\eta-\Sigma(\boldsymbol{k},\omega+\mathrm{i}\eta)}
\end{equation}
对于自由电子系统，由于不存在相互作用 $H_{\mathrm{int}}=0$，自然应当有 $\Sigma^{(0)}(\boldsymbol{k},\omega+\mathrm{i}\eta)=0$，因此自由电子的推迟格林函数为
\begin{equation}
G_{\mathrm{r}}^{(0)}(\boldsymbol{k},\omega)
=\frac{1}{\omega-E_{k}+\mathrm{i}\eta}
\end{equation}
以后将推迟格林函数的严格自能记为
\begin{equation}
\Sigma_{\mathrm{r}}(\boldsymbol{k},\omega)
\equiv\Sigma(\boldsymbol{k},\omega+\mathrm{i}\eta)
=\mathrm{Re}\Sigma_{\mathrm{r}}
+\mathrm{i}\,\mathrm{Im}\Sigma_{\mathrm{r}}
\end{equation}
一般而言，自能是 $\omega$ 的复函数，这时利用(C.16)可将相互作用系统的格林函数，式(C.15)，改写成
\begin{equation}
G_{\mathrm{r}}(\boldsymbol{k},\omega)
=\frac{1}{[G_{\mathrm{r}}^{(0)}(\boldsymbol{k},\omega)]^{-1}
-\Sigma_{\mathrm{r}}(\boldsymbol{k},\omega)}
\end{equation}
此即著名的戴逊方程. 在量子多体的专著中有时也表示为下列两种等效的形式：
\begin{equation}
[G_{\mathrm{r}}^{(0)}(\boldsymbol{k},\omega)]^{-1}
-[G_{\mathrm{r}}(\boldsymbol{k},\omega)]^{-1}
=\Sigma_{\mathrm{r}}(\boldsymbol{k},\omega)
\end{equation}
或
\begin{equation}
G_{\mathrm{r}}(\boldsymbol{k},\omega)
=G_{\mathrm{r}}^{(0)}(\boldsymbol{k},\omega)
+G_{\mathrm{r}}^{(0)}(\boldsymbol{k},\omega)
\Sigma_{\mathrm{r}}(\boldsymbol{k},\omega)
G_{\mathrm{r}}(\boldsymbol{k},\omega)
\end{equation}
由于自能函数 $\Sigma_{\mathrm{r}}(\boldsymbol{k},\omega)$ 本身含有限的虚部，因此式(C.18)右边 $[G_{\mathrm{r}}^{(0)}(\boldsymbol{k},\omega)]^{-1}$ 中的 $\mathrm{i}\eta$ 可略去不计，而将戴逊方程简单表示为
\begin{equation}
G_{\mathrm{r}}(\boldsymbol{k},\omega)
=[\omega-E_{k}-\Sigma_{\mathrm{r}}(\boldsymbol{k},\omega)]^{-1}
\end{equation}
具体计算自能函数必须作近似，其中最为有效的是图形微扰方法（包括部分求和法）. 应用戴逊方程分析费密面附近准粒子的性质和论证相互作用电子系统中存在费密面都将是十分方便的，具体的讨论可参阅第四章§11.
